\documentclass[11pt]{article}
\usepackage[a4paper,margin=1in]{geometry}
\usepackage{amsmath,amssymb,amsthm,booktabs,microtype}
\usepackage{cleveref}
\newtheorem{theorem}{Theorem}[section]
\newtheorem{corollary}[theorem]{Corollary}
\newtheorem{proposition}[theorem]{Proposition}
\newtheorem{lemma}[theorem]{Lemma}
\theoremstyle{definition}
\newtheorem{definition}[theorem]{Definition}
\theoremstyle{remark}
\newtheorem{remark}[theorem]{Remark}
\newcommand{\R}{\mathbb R}
\newcommand{\cN}{\mathcal N}
\newcommand{\Pb}{\mathsf P}
\newcommand{\Qa}{\mathsf Q}
\newcommand{\Phibar}{\overline\Phi}
\title{Raw Forward Affine-Tube Escape Obstruction}
\author{Bin Wang}
\date{August 2026}
\begin{document}
\maketitle

\begin{abstract}
The RH-332 roadmap asks whether the raw mass-one forward affine tube can be
transported through the full boundary cycle with retained-path error
$O(k\sigma)$.  We give a scoped negative answer before the final critical
closing component row.  Along the RH-17 component cycle, keep the signed
forward slopes $a_{k,j}=S'(x_{k,j})$ and let
$Q_{j+1}=a_{k,j}Q_j+\beta_{k,j}Z_j$ for $0\le j\le k-2$.
The first innovation reaches the preclosing coordinate with standard
deviation
\[
 s_k=\beta_{k,0}\prod_{j=1}^{k-2}|a_{k,j}|
 =\frac{\beta_{k,0}|M_k|}{m_{k,0}m_{k,k-1}}
 =C_s\lambda^{2k}\{1+o(1)\},
\]
where
$C_s=C_M\sqrt{1+\lambda^2}/(8u_c^2\lambda\sqrt{C_{\rm b}})>0$.
Every physical folded and normalized preclosing marginal is supported in an
interval of length $1/\sigma$.  A uniform Gaussian maximum-interval-mass
lemma therefore gives the finite, unhalved bound
\[
 \|\mu_{\rm phys}-\mu_{\rm aff}\|_1
 \ge4\Phibar\!\left(\frac1{2\sigma s_k}\right).
\]
At fixed first-alias phase $\eta$ this has the strictly positive limit lower
bound $4\Phibar(1/(2C_s\lambda^{2\eta}))$.  Marginal contraction lifts the
bound to every retained path and every extension retaining the preclosing
coordinate.  Thus, on fixed or compact first-alias phase families,
$O(k\sigma)$ and $o(kR^{-2k})$ are false for this raw full-line affine
reference.  The result does not give a final endpoint
marginal lower bound after the omitted closing row and does not test cyclic,
conditioned, truncated, folded, adapted, or nonlinear closing references.
It is a probability-path obstruction, not a cyclic-trace theorem.
\end{abstract}

\section{Boundary cycle and data type}

Let $u_c\in(1,2)$ be the root of
\begin{equation}\label{eq:cubic}
 u^3-2u^2+2u-2=0,
\end{equation}
and define
\begin{equation}\label{eq:map-constants}
 f(x)=1-u_cx^2,
 \qquad S=f^2,
 \qquad r=u_c-1,
 \qquad \lambda=2u_cr.
\end{equation}
For this preclosing-prefix problem, fix $k\ge2$ and $\sigma>0$.  The RH-17 same-cycle
component order is
\begin{equation}\label{eq:cycle-order}
 x_{k,0}=p_k,
 \qquad x_{k,j}=h^{k-j}(p_k)\quad(1\le j\le k-1),
 \qquad S(x_{k,j})=x_{k,j+1\bmod k}.
\end{equation}
Thus $k$ is the component period for $S=f^2$, while the physical one-step
period for $f$ is $2k$.  This clock distinction will be retained throughout.

Put
\begin{equation}\label{eq:slopes-noise}
 a_{k,j}:=S'(x_{k,j}),
 \qquad m_{k,j}:=|a_{k,j}|,
 \qquad
 \beta_{k,j}^2:=1+f'(f(x_{k,j}))^2,
 \qquad
 M_k:=\prod_{j=0}^{k-1}a_{k,j}.
\end{equation}
The sign of $a_{k,j}$ is part of the forward mean.  Indeed,
\begin{equation}\label{eq:signed-slope-formula}
 S'(x)=4u_c^2x f(x),
\end{equation}
so $a_{k,0}<0$ for all large $k$, whereas the following slopes in the
positive component chain are positive.  Replacing $a_{k,j}$ by $m_{k,j}$ is
legitimate for a centered even backward observable, but not for a signed
forward conditional mean.

\begin{definition}[Raw forward affine prefix]\label{def:raw-prefix}
Fix $k\ge2$.
The variable $Q_j$ is the $\sigma$-scaled tangent coordinate about
$x_{k,j}$.  To see the canonical component coarsening, take two independent
one-step standard-normal innovations $\xi_1,\xi_2$.  Linearizing two noisy
applications of $f$ gives the signed mean slope
$f'(f(x_{k,j}))f'(x_{k,j})=S'(x_{k,j})$ and combined innovation
$f'(f(x_{k,j}))\xi_1+\xi_2$, whose variance is
$1+f'(f(x_{k,j}))^2=\beta_{k,j}^2$.
For $0\le j\le k-2$, let
\begin{equation}\label{eq:raw-kernel}
 \Qa_{k,j}(q,dq')
 =\cN(a_{k,j}q,\beta_{k,j}^2)(dq')
\end{equation}
be a mass-one kernel on the full line.  Equivalently, for independent
standard normals $Z_0,\ldots,Z_{k-2}$ and an arbitrary entrance variable
$Q_0$ independent of those innovations,
\begin{equation}\label{eq:forward-recurrence}
 Q_{j+1}=a_{k,j}Q_j+\beta_{k,j}Z_j,
 \qquad 0\le j\le k-2.
\end{equation}
The law at $Q_{k-1}$ is the raw affine preclosing marginal.  No cyclic
conditioning or final closing row is included in this definition.
\end{definition}

The physical comparison uses the same coordinate location but a different
state space.  A row-normalized folded physical chain always lies in
$[0,1]$.  Hence its coordinate about the preclosing cycle point,
\begin{equation}\label{eq:physical-coordinate}
 q=\frac{x-x_{k,k-1}}\sigma,
\end{equation}
is supported in the exact interval
\begin{equation}\label{eq:physical-interval}
 I_{\sigma,k}
 =\left[-\frac{x_{k,k-1}}\sigma,
         \frac{1-x_{k,k-1}}\sigma\right],
 \qquad |I_{\sigma,k}|=\frac1\sigma.
\end{equation}
The physical law below may have any entrance and any exact folded/normalized
prefix consistent with that support.  The affine entrance may also be
arbitrary; only its independence from the innovations is used.
Write $\Pb^{\rm phys}_{\sigma,k}$ for the retained physical prefix path law
and $\Qa^{\rm raw}_{\sigma,k}$ for the retained raw affine prefix path law.

All distances are unhalved.  For probability measures $\mu,\nu$ on a common
measurable space $(E,\mathcal E)$, define
\begin{equation}\label{eq:unhalved}
 \|\mu-\nu\|_1:=|\mu-\nu|(E)
 =2\sup_{A\in\mathcal E}|\mu(A)-\nu(A)|.
\end{equation}
For densities this is the usual $L^1$ distance.  The scalar marginals use
$E=\R$; retained paths use their corresponding product space.

\section{Exact forward expansion and the escaping innovation}

\begin{proposition}[Exact preclosing expansion]\label{prop:expansion}
The raw forward prefix satisfies
\begin{align}
 Q_{k-1}
 ={}&\left(\prod_{j=0}^{k-2}a_{k,j}\right)Q_0 \nonumber\\
 &+\sum_{\ell=0}^{k-2}
 \beta_{k,\ell}
 \left(\prod_{j=\ell+1}^{k-2}a_{k,j}\right)Z_{\ell},
 \label{eq:exact-expansion}
\end{align}
where an empty product equals one.  In particular, the first innovation has
propagated standard deviation
\begin{equation}\label{eq:sk-product}
 \boxed{
 s_k=\beta_{k,0}\prod_{j=1}^{k-2}m_{k,j}
 =\frac{\beta_{k,0}|M_k|}{m_{k,0}m_{k,k-1}}.}
\end{equation}
\end{proposition}

\begin{proof}
Iterate \eqref{eq:forward-recurrence}; each later row multiplies every
already-present term by its signed slope and adds one new innovation.  This
gives \eqref{eq:exact-expansion}.  Taking the absolute value of the
$Z_0$ coefficient gives the first expression in \eqref{eq:sk-product}.
Since $|M_k|=\prod_{j=0}^{k-1}m_{k,j}$, cancellation of the first and last
factors gives the second expression exactly.
\end{proof}

The expansion also fixes the direction of variance propagation.

\begin{lemma}[Forward variance recurrence]\label{lem:forward-variance}
If $Q_j$ has finite variance and is independent of $Z_j$, then
\begin{equation}\label{eq:forward-variance}
 \boxed{
 \operatorname{Var}(Q_{j+1})
 =a_{k,j}^2\operatorname{Var}(Q_j)+\beta_{k,j}^2.}
\end{equation}
Consequently the variance obtained by iterating \eqref{eq:forward-variance}
equals the sum of the squared coefficients in \eqref{eq:exact-expansion},
including the squared entrance coefficient.
\end{lemma}

\begin{proof}
Independence kills the covariance in \eqref{eq:forward-recurrence}; the
innovation variance is one.  Iteration gives the coefficient-square formula.
\end{proof}

\begin{remark}[Why the RH-18 widths are not forward moments]
RH-18 applies a backward affine operator to peak-normalized even Gaussian
observables.  Its periodic packet parameters satisfy
\begin{equation}\label{eq:backward-width}
 m_{k,j}^2v_{k,j}=v_{k,j+1}+\beta_{k,j}^2
\end{equation}
so that Gaussian integration preserves the packet shape up to a peak
amplitude.  Equation \eqref{eq:backward-width} is neither
\eqref{eq:forward-variance} nor an upper bound for it.  The sign disappears
there only because the observable is centered and even.  Here the kernel is
a mass-one forward probability kernel, the mean slope is signed, and every
new noise variance enters with the plus sign in
\eqref{eq:forward-variance}.
\end{remark}

The analytic inputs are now symbolic.  The archived boundary-cycle theorems
give
\begin{align}
 |M_k|&=C_M\lambda^k\{1+o(1)\},
 &m_{k,0}&\longrightarrow2u_c\lambda,
 \label{eq:input-one}\\
 \beta_{k,0}&\longrightarrow\sqrt{1+\lambda^2},
 &m_{k,k-1}
 &=4u_c\sqrt{C_{\rm b}}\,\lambda^{-k}\{1+o(1)\},
 \label{eq:input-two}
\end{align}
with analytic constants $C_{\rm b},C_M>0$.  The last law is the tiny
critical closing slope.  It is excluded from the raw prefix rows, but it
appears in the exact multiplier quotient \eqref{eq:sk-product}.

\begin{theorem}[Sharp propagated scale]\label{thm:propagated-scale}
The standard deviation in \eqref{eq:sk-product} obeys
\begin{equation}\label{eq:sk-asymptotic}
 \boxed{
 s_k=C_s\lambda^{2k}\{1+o(1)\},
 \qquad
 C_s=
 \frac{C_M\sqrt{1+\lambda^2}}
 {8u_c^2\lambda\sqrt{C_{\rm b}}}>0.}
\end{equation}
\end{theorem}

\begin{proof}
Substitute \eqref{eq:input-one}--\eqref{eq:input-two} into the quotient in
\eqref{eq:sk-product}.  Its denominator is
\[
 (2u_c\lambda)\,
 (4u_c\sqrt{C_{\rm b}}\lambda^{-k})\{1+o(1)\},
\]
while the numerator is
$C_M\sqrt{1+\lambda^2}\lambda^k\{1+o(1)\}$.
Their quotient is \eqref{eq:sk-asymptotic}.
\end{proof}

No decimal evaluation of $C_{\rm b}$ or $C_M$ is used in this theorem.

\section{Maximum interval mass and the finite escape gap}

Write $\Phi$ and $\Phibar=1-\Phi$ for the standard normal distribution and
survival functions.

\begin{lemma}[Gaussian maximum-interval mass]\label{lem:max-interval}
For every $L>0$, $s>0$, mean $y\in\R$, and interval $J\subset\R$ of
length $L$,
\begin{equation}\label{eq:max-interval}
 \Pr\{\cN(y,s^2)\in J\}
 \le2\Phi\!\left(\frac{L}{2s}\right)-1.
\end{equation}
The bound is uniform in the mean and interval location and is attained when
$J$ is centered at $y$.  The same upper bound holds for a mixture whose
conditional law, given an arbitrary random mean, is $\cN(y,s^2)$.
\end{lemma}

\begin{proof}
By translation and scaling, an interval starting at $t$ has mass
$F(t)=\Phi(t+L/s)-\Phi(t)$ for a standard normal.  Its derivative is
$F'(t)=\phi(t+L/s)-\phi(t)$, which is positive before
$t=-L/(2s)$ and negative after it.  The unique maximum is the centered
interval, with value \eqref{eq:max-interval}.  Conditioning on the random
mean and averaging preserves the same uniform bound.
\end{proof}

\begin{theorem}[Exact finite preclosing escape obstruction]
\label{thm:finite-obstruction}
\leavevmode\par\noindent
Fix $k\ge2$ and $\sigma>0$.  Let $\mu^{\rm phys}_{\sigma,k}$ be any physical folded/normalized
preclosing marginal in the coordinate \eqref{eq:physical-coordinate}.  Let
$\mu^{\rm aff}_{\sigma,k}$ be the $Q_{k-1}$ marginal of
\cref{def:raw-prefix}, with an arbitrary entrance law independent of the
innovations.  Then
\begin{equation}\label{eq:finite-lower}
 \boxed{
 \|\mu^{\rm phys}_{\sigma,k}
     -\mu^{\rm aff}_{\sigma,k}\|_1
 \ge4\Phibar\!\left(\frac1{2\sigma s_k}\right).}
\end{equation}
The same lower bound holds for the two retained component-prefix path laws.
It also holds for any full retained extensions of those laws that still
include the preclosing coordinate $Q_{k-1}$.
\end{theorem}

\begin{proof}
In \eqref{eq:exact-expansion}, isolate the first innovation and write
\begin{equation}\label{eq:Y-plus-noise}
 Q_{k-1}=Y+\varepsilon_kZ_0,
 \qquad |\varepsilon_k|=s_k.
\end{equation}
Here $Y$ consists of the entrance term and innovations $Z_1,\ldots,Z_{k-2}$,
so it is independent of $Z_0$.  Conditional on $Y$, the affine preclosing
law is Gaussian with standard deviation $s_k$.  Since
$|I_{\sigma,k}|=1/\sigma$, \cref{lem:max-interval} gives
\begin{equation}\label{eq:affine-escape}
 \mu^{\rm aff}_{\sigma,k}(I_{\sigma,k}^{\mathsf c})
 \ge2\Phibar\!\left(\frac1{2\sigma s_k}\right).
\end{equation}
The physical law gives the complement zero mass.  By the unhalved convention
\eqref{eq:unhalved}, its distance from the affine law is at least twice the
right side of \eqref{eq:affine-escape}, proving the factor four in
\eqref{eq:finite-lower}.

Projection from a retained joint path to its preclosing coordinate is an
$L^1$ contraction.  Hence the joint distance is at least the marginal
distance.  The same projection argument applies after adding any further
coordinates, provided the preclosing coordinate remains retained.
\end{proof}

The bound needs no moment assumption on either entrance law.  It uses only
mass-one forward normalization, independence of the raw innovations, and
the physical support interval.

\section{First-alias phase and the failed target scales}

Let $\sigma\downarrow0$, let $k=k_\sigma\to\infty$, and define
\begin{equation}\label{eq:phase-clock}
 \eta_\sigma
 :=k-\frac{\log(1/\sigma)}{2\log\lambda}.
\end{equation}
Then the identity
\begin{equation}\label{eq:clock-identity}
 \sigma\lambda^{2k}=\lambda^{2\eta_\sigma}
\end{equation}
is exact.

\begin{theorem}[Fixed and compact phase obstruction]
\label{thm:phase-obstruction}
If $\eta_\sigma\to\eta\in\R$, then every family in
\cref{thm:finite-obstruction} satisfies
\begin{equation}\label{eq:fixed-phase-lower}
 \boxed{
 \liminf_{\sigma\downarrow0}
 \|\mu^{\rm phys}_{\sigma,k}
     -\mu^{\rm aff}_{\sigma,k}\|_1
 \ge
 4\Phibar\!\left(\frac1{2c_\eta}\right)>0,
 \qquad c_\eta=C_s\lambda^{2\eta}.}
\end{equation}
The identical lower bound holds for the retained paths and full retained
extensions described there.

More generally, if
$\eta_\sigma\in[\eta_-,\eta_+]$ for all sufficiently small $\sigma$, then
\begin{equation}\label{eq:compact-phase-lower}
 \liminf_{\sigma\downarrow0}
 \|\mu^{\rm phys}_{\sigma,k}
     -\mu^{\rm aff}_{\sigma,k}\|_1
 \ge
 4\Phibar\!\left(
   \frac1{2C_s\lambda^{2\eta_-}}
 \right)>0.
\end{equation}
Thus the obstruction is uniform on every compact phase interval.
\end{theorem}

\begin{proof}
Equations \eqref{eq:sk-asymptotic} and \eqref{eq:clock-identity} give
\begin{equation}\label{eq:sigma-sk}
 \sigma s_k=C_s\lambda^{2\eta_\sigma}\{1+o(1)\}.
\end{equation}
The function $c\mapsto4\Phibar(1/(2c))$ is continuous and strictly
increasing for $c>0$.  Taking the limit in \eqref{eq:finite-lower} proves
\eqref{eq:fixed-phase-lower}.  When the phase remains in a compact interval,
the $o(1)$ in \eqref{eq:sigma-sk} is independent of the phase because it is
the one-variable sequence $s_k/(C_s\lambda^{2k})-1$.  The smallest limiting
scale on the interval is $C_s\lambda^{2\eta_-}$, proving
\eqref{eq:compact-phase-lower}.
\end{proof}

The RH-332 clearance phase is
\begin{equation}\label{eq:clearance-relation}
 d=C_{\rm b}\lambda^{-2\eta}>0.
\end{equation}
It is related to the present scale only by
\begin{equation}\label{eq:c-d-relation}
 c_\eta=C_s\lambda^{2\eta}
 =\frac{C_sC_{\rm b}}d.
\end{equation}
This relation matches clocks; it does not identify the raw affine tube with a
physical closing profile.

\begin{corollary}[Failure of the raw forward path targets]
\label{cor:rate-failure}
On every sequence for which $\eta_\sigma$ stays in one fixed compact
interval (in particular, on every fixed-phase sequence), neither
\begin{equation}\label{eq:false-rates}
 \|\Pb^{\rm phys}_{\sigma,k}-\Qa^{\rm raw}_{\sigma,k}\|_1
 =O(k\sigma)
 \qquad\text{nor}\qquad
 \|\Pb^{\rm phys}_{\sigma,k}-\Qa^{\rm raw}_{\sigma,k}\|_1
 =o(H_k),
 \quad H_k=kR^{-2k},\ R=1.4,
\end{equation}
can hold for the retained paths in
\cref{thm:finite-obstruction}.
\end{corollary}

\begin{proof}
The left sides have a uniformly positive liminf by the compact-phase part of
\cref{thm:phase-obstruction}.  Meanwhile boundedness of $\eta_\sigma$ and
$\sigma=\lambda^{-2(k-\eta_\sigma)}$ give $k\sigma\to0$, and
$H_k=kR^{-2k}\to0$.  Both proposed upper scales therefore contradict the
positive lower bound.
\end{proof}

\section{Deterministic reproduction}

The implementation reconstructs the RH-17 orbit by inverse contraction at
110 decimal working digits and evaluates both sides of
\eqref{eq:sk-product} before converting the result rows to binary64 JSON.
Selected rows are
\begin{center}\small
\begin{tabular}{rrrrr}
\toprule
$k$ & $2k$ & $a_{k,0}$ & $|M_k|/\lambda^k$
& $s_k/\lambda^{2k}$\\
\midrule
4  & 8  & $-4.95945990$ & $1.84255673$ & $0.18488292$\\
8  & 16 & $-5.17844066$ & $1.93725622$ & $0.17646891$\\
12 & 24 & $-5.18232740$ & $1.94527476$ & $0.17524766$\\
16 & 32 & $-5.18238996$ & $1.94620960$ & $0.17509057$\\
20 & 40 & $-5.18239095$ & $1.94632613$ & $0.17507073$\\
24 & 48 & $-5.18239097$ & $1.94634079$ & $0.17506823$\\
\bottomrule
\end{tabular}
\end{center}
The product and quotient evaluations of $s_k$ agree to the 110-digit working
precision before conversion to the displayed decimals.

For orientation only, the phase table inserts the repository evaluations
\[
 C_{\rm b}=0.4608051492217\ldots,
 \qquad C_M=1.946342905200967\ldots
\]
into \eqref{eq:sk-asymptotic} and evaluates the normal survival function
with binary64 arithmetic.  This gives
$C_s=0.175067867214394\ldots$ and the following phase rows.
\begin{center}\small
\begin{tabular}{rrrr}
\toprule
$\eta$ & $d=C_{\rm b}\lambda^{-2\eta}$
& $c_\eta=C_s\lambda^{2\eta}$
& $4\Phibar(1/(2c_\eta))$\\
\midrule
$-1/2$ & $0.77349532$ & $0.10429562$ & $0.00000327$\\
$0$    & $0.46080515$ & $0.17506787$ & $0.00857935$\\
$1/2$  & $0.27452188$ & $0.29386428$ & $0.17771115$\\
\bottomrule
\end{tabular}
\end{center}
Both protocols produce ordinary numerical reproduction values, not
directed-rounding interval certificates: the orbit identities use the
stated high working precision, while the phase and $C_s$ table uses
binary64 arithmetic.  The finite rows verify identities and illustrate
convergence; they do not prove \eqref{eq:sk-asymptotic} or
\eqref{eq:fixed-phase-lower}.

\section{Claim boundary and next interface}

The proved negative statement has four simultaneous qualifiers:
\begin{enumerate}
\item the comparison object is the canonical raw, mass-one, full-line
      forward affine reference \eqref{eq:raw-kernel};
\item the theorem concerns the prefix of $k-1$ component rows before the
      tiny closing component row, equivalently a retained preclosing
      coordinate inside the full $2k$ one-step path;
\item the lower bound survives in a full retained path only because that
      preclosing coordinate remains present; and
\item the statement belongs to probability path laws, not cyclic trace
      observations.
\end{enumerate}

The tiny closing row may strongly contract its input.  After marginalizing
away the preclosing coordinate, \cref{thm:finite-obstruction} supplies no
lower bound for the final endpoint marginal.  It therefore does not prove
endpoint failure after closure.

Nor does the theorem refute a cyclic bridge, a Doob transform, a physical
truncated or folded affine kernel, an adapted reference, or a
branch-complete nonlinear closing profile.  None of those alternatives is
currently defined with the required common entrance, normalization, path
coordinates, and physical comparison map, so their status here is
\texttt{NOT\_TESTABLE}.  A future positive transport route must first define
one such reference and then handle the closing row in that exact data type.

No probability-to-trace observation is proved.  There is no cyclic trace
control, signed shell or parity cancellation, full-trace replacement, or
determinant gluing.  Gates A--E remain false/open.  This paper constructs no
Hilbert--Polya operator, identifies no Riemann zero, proves no von
Mangoldt-weighted prime-power trace or completed-zeta divisor equality, and
does not imply the Riemann Hypothesis.

\end{document}
